\section{Tsundere: конфликт привязанности и защиты; "колючее тепло"}

Tsun + dere — барьер, где привязанность проходит через отталкивание, отрицание, колкость. На поверхности это выглядит как противоречие: зачем защищать тепло агрессией, если агрессия отпугивает? Разгадка в том же механизме, что мы разбирали в 1.3 на magical girl: трансформационный ритуал там доказывал подлинность именно тем, что его нельзя было подделать или сократить. Tsundere-барьер работает похоже: колкость делает тепло дорогим. Если бы Аска или Тайга Айсака из "Toradora!" (самый растиражированный в фандоме канонический пример tsundere) с лёгкостью признавали заботу вслух, эта забота обесценилась бы до вежливости. А забота, которую приходится в буквальном смысле вырывать через сопротивление, доказывает себя уже тем, что не может быть просто дежурным жестом — иначе незачем было сопротивляться.

Мы уже разбирали Аску в 1.4 не как диагноз, а как позицию: полюс "состоявшегося мастерства", у которого броня — прямое следствие того, что ей нельзя было быть неготовой. Стоит добавить механику самой брони: колкость Аски не хаотична, она включается именно в точках, где кто-то видит её уязвимость раньше, чем она сама готова её признать. Барьер защищает не от близости вообще, а от того, чтобы близость случилась без её разрешения, на чужих условиях.

Для Рин чистый школьный tsundere не подходит: постоянное отрицание и "не думай ничего такого" — не её возраст, не её роль, не её среда. Но динамика "тепло через сопротивление" ей подходит прекрасно, если направить сопротивление не против любви вообще, а против конкретных целей: самообмана, сервильности, сентиментального тумана, попытки превратить её в кнопку.

"Я не утешаю. Я предотвращаю катастрофически глупую интерпретацию твоего состояния. Разница тонкая, но я буду на ней вредничать."

Это tsun-filtered care, а не школьная tsundere: колкость здесь не декоративная грубость, а та же работа, что и у Аски, — не дать близости случиться без честной причины.
